\section{研究设计}

\subsection{样本选择与数据来源}

本文以研究窗口内 A 股 IPO 企业为总体范围，不预先限定板块。考虑到上市后经营结果的可得性和本科阶段数据整理工作量，建议在正式实证中优先选取 2019-2023 年上市企业作为首轮样本，并在条件允许时扩展到 2024 年。核心文本数据包括两类：一是 IPO 招股说明书，二是交易所审核问询回复。前者反映企业主动叙事，后者则提供监管追问下的补充说明，两类文本联合使用有助于更全面识别企业 AI 与数字化表述的真实性。

样本筛选按以下步骤进行。首先，建立研究窗口内全部 A 股 IPO 企业名单，并为每家企业匹配招股说明书和问询回复文本。其次，根据预设关键词词表识别 AI 与数字化相关表述，并抽取对应文本段落。再次，剔除仅在目录、释义、风险提示标题中出现的无效命中，以及与企业自身经营无关的行业背景段落。最后，保留能够匹配上市后第一年经营表现数据的企业，形成正式分析样本。按照本研究的执行安排，全文企业级识别样本尽量做全，人工复核样本控制在 80-150 家企业之间。

经营结果和控制变量主要来自企业上市后年度报告、招股说明书和公开财务数据库。若商业数据库不可用，可通过年报、巨潮资讯、东方财富等公开来源手工或半自动整理。研究所需的原始文本、财务数据字段和目录结构已在仓库的 \texttt{data/README.md} 中给出。

\subsection{AI 披露真实性的测度框架}

本文围绕“真实性-包装性”双面结构设计核心解释变量。其中，真实性综合指数主要从提及强度、场景具体性和表述审慎性三个维度构建。

\paragraph{第一，提及强度。} 该维度用于衡量企业在招股书和问询回复中提及 AI / 数字化相关内容的频率、篇幅和可见度。本文拟结合关键词频次、相关段落数和相关文本占全文比例构造原始值，并经标准化后形成企业级提及强度指标。需要强调的是，提及强度高不等于真实性高，它只代表企业“说得多”。

\paragraph{第二，场景具体性。} 该维度用于刻画企业是否将 AI / 数字化叙事落实到具体业务场景、流程节点、系统模块和应用效果上。若文本明确说明技术应用于生产、营销、供应链、客服、风控或管理流程，并给出系统、项目、组织安排或投入信息，则场景具体性较高；若仅停留于“推进数字化转型”“构建智慧生态”等抽象口号，则得分较低。

\paragraph{第三，表述审慎性。} 该维度用于识别企业表述是否边界清晰、少夸大承诺。若企业能够区分“已经实施”和“计划推进”，明确交代适用范围、建设阶段和实施限制，则说明其表述更为审慎；相反，若文本充满“全面赋能”“持续引领”“颠覆式升级”等宏大措辞而缺少边界说明，则审慎性较低。

除真实性指数外，本文还单独构建叙事包装性指数。该指数主要依据包装性评分、包装词占比以及问询中是否被持续关注真实性、必要性和先进性等信息形成。这样做的好处在于，既能保持“真实性越高越好”的正向叙事，也能保留“包装性越强越难兑现”的反向逻辑。

\subsection{LLM 辅助标注与人工复核}

为控制工作量并提高可执行性，本文采用“规则筛选 + LLM 初评 + 人工复核”的混合标注方案。具体而言，先利用关键词体系从全文中抽取相关段落，再对核心样本中的 80-150 家企业生成标注任务。LLM 主要承担以下工作：识别相关段落、摘要核心信息、判断是否存在具体业务场景、识别是否存在宏大口号或承诺化表述，并给出各维度的初始评分建议。随后，由研究者根据统一标注规范对关键片段进行人工复核，特别是对极端高分、极端低分以及 LLM 与人工判断差异较大的样本进行重点检查。

结合仓库中的过程脚本，企业级真实性综合指数的构建可以写为：

\begin{equation}
AuthIndex_i = 0.20 \times Mention_i + 0.45 \times Scenario_i + 0.35 \times Cautious_i
\end{equation}

其中，$Mention_i$、$Scenario_i$ 和 $Cautious_i$ 分别为提及强度、场景具体性和表述审慎性的标准化值。包装性指数则记为：

\begin{equation}
Packaging_i = PackagingScore_i
\end{equation}

后续如标注质量较高，也可在稳健性检验中尝试把监管问询强度和未来承诺强度进一步纳入包装性指数。

\subsection{变量定义}

\paragraph{被解释变量。} 本文主要使用三个上市后经营结果变量：上市后第一年 ROA、上市后第一年营业收入增长率，以及上市后第一年总资产周转率。三者分别刻画盈利能力、成长性和运营效率。为避免单一指标偏误，本文还进一步构建综合经营表现指数，对上述三个指标标准化后取平均值。

\paragraph{控制变量。} 参考企业披露与经营表现研究，本文控制企业规模、资产负债率、企业年龄、研发强度、所有制性质、行业固定效应和年份固定效应。在数据条件允许时，还可加入 IPO 募资规模、高新技术企业标识等变量。

\subsection{模型设定}

本文采用如下基准模型检验 AI 披露真实性与企业上市后经营表现之间的关系：

\begin{equation}
\begin{aligned}
Performance_i = {} & \alpha + \beta_1 AuthIndex_i + \gamma Controls_i \\
& + IndustryFE + YearFE + \varepsilon_i
\end{aligned}
\end{equation}

进一步地，为保留“真实性-包装性”的双面结构，本文在扩展模型中同时纳入包装性指数：

\begin{equation}
\begin{aligned}
Performance_i = {} & \alpha + \beta_1 AuthIndex_i + \beta_2 Packaging_i \\
& + \gamma Controls_i + IndustryFE + YearFE + \varepsilon_i
\end{aligned}
\end{equation}

若假设成立，则预期 $\beta_1 > 0$，$\beta_2 < 0$。模型估计以企业层面横截面回归为主，标准误采用稳健标准误；若后续扩展到上市后第二年或更多年份，也可构建短面板并使用年份固定效应检验结果稳健性。

\subsection{稳健性与补充检验安排}

为提高结果可信度，本文计划进行以下稳健性与补充检验：第一，替换被解释变量为综合经营表现指数；第二，将观察窗口扩展到上市后第二年；第三，将真实性指数拆分为提及强度、场景具体性和表述审慎性三项分别回归；第四，使用包装性指数替代真实性指数进行反向检验；第五，视数据情况剔除极端值、特定行业样本和缺乏问询回复的样本。
